\documentclass[11pt]{article}

\usepackage[margin=1in]{geometry}
\usepackage[T1]{fontenc}
\usepackage{lmodern}
\usepackage{amsmath}
\usepackage{booktabs}
\usepackage{enumitem}
\usepackage{hyperref}
\usepackage{xcolor}
\usepackage{listings}

\hypersetup{
    colorlinks=true,
    linkcolor=blue,
    urlcolor=blue,
    citecolor=blue
}

\lstset{
    basicstyle=\ttfamily\small,
    breaklines=true,
    frame=single,
    columns=fullflexible,
    showstringspaces=false
}
\title{ADR-001\\Compile-Time Policy Architecture}
\author{William Gordon Carter}
\date{\today}

\begin{document}

\maketitle

\begin{center}
\begin{tabular}{ll}
\toprule
\textbf{Field} & \textbf{Value} \\
\midrule
ADR & 001 \\
Title & Compile-Time Policy Architecture \\
Status & Proposed \\
Author & William Gordon Carter \\
Created & \today \\
Supersedes & None \\
Superseded By & --- \\
\bottomrule
\end{tabular}
\end{center}

\tableofcontents
\newpage

\section{Context}

NavKit is intended to be a modular estimation and navigation development
platform for embedded aerospace systems. Many of the most important choices in
such a system are not runtime choices. They are compile-time engineering
assumptions: state layout, measurement structure, process model, injection
strategy, reset strategy, planet, gravity, mechanization, frame convention, and
simulation environment.

This ADR defines the broad architectural pattern NavKit uses to represent
those assumptions.

\section{Decision}

NavKit shall use a compile-time policy architecture based on:

\begin{center}
\textbf{Concept} $\rightarrow$ \textbf{CRTP Policy Base} $\rightarrow$
\textbf{Concrete Policy} $\rightarrow$ \textbf{Runtime Algorithm}
\end{center}

Concepts define compile-time capabilities. CRTP policy bases provide shared
static implementation. Concrete policies aggregate the facts, equations, or
behavior needed for a particular configuration. Runtime algorithms consume
policies through constrained template parameters.

\section{Core Idea}

A C++20 concept is not a base class and is not explicitly implemented by a
concrete type. A concept is a compile-time predicate. It asks whether a type
satisfies a set of requirements.

For example:

\begin{lstlisting}[language=C++]
template<typename T>
concept PlanetPolicy = requires
{
    T::mu_m3_s2;
    typename T::InertialFrame;
    typename T::FixedFrame;
};
\end{lstlisting}

A concrete type does not declare that it implements \texttt{PlanetPolicy}. It
simply provides the required members:

\begin{lstlisting}[language=C++]
struct Wgs84
{
    using InertialFrame = frames::EarthCenteredInertial;
    using FixedFrame = frames::EarthCenteredEarthFixed;

    static inline constexpr double mu_m3_s2 = 3.986004418e14;
};
\end{lstlisting}

The concept becomes useful when it constrains a template:

\begin{lstlisting}[language=C++]
template<PlanetPolicy Planet>
struct SphericalGravity;
\end{lstlisting}

The syntax above is shorthand for:

\begin{lstlisting}[language=C++]
template<typename Planet>
requires PlanetPolicy<Planet>
struct SphericalGravity;
\end{lstlisting}

At this point, the compiler checks whether the supplied template argument
satisfies \texttt{PlanetPolicy}.

\section{Terminology}

\subsection{Policy}

A policy is a compile-time type that defines an engineering choice. Policies
may define constants, equations, behavior, static helpers, frame types, or
compile-time metadata.

Examples include:

\begin{itemize}
    \item \texttt{PlanetPolicy}
    \item \texttt{GravityPolicy}
    \item \texttt{PropagationPolicy}
    \item \texttt{MechanizationPolicy}
    \item \texttt{MeasurementPolicy}
    \item \texttt{InjectionPolicy}
    \item \texttt{ResetPolicy}
    \item \texttt{NoisePolicy}
\end{itemize}

\subsection{Concept}

A concept defines compile-time requirements. In NavKit, concepts should usually
represent capabilities rather than inheritance hierarchies.

Examples:

\begin{itemize}
    \item \texttt{PlanetPolicy}
    \item \texttt{RotatingPlanetPolicy}
    \item \texttt{EllipsoidPlanetPolicy}
    \item \texttt{J2PlanetPolicy}
\end{itemize}

A single concrete policy may satisfy many concepts without explicitly declaring
any of them.

\subsection{CRTP Policy Base}

A CRTP policy base provides shared static functionality for a family of
policies.

Example:

\begin{lstlisting}[language=C++]
template<typename Derived>
struct PlanetPolicyBase
{
    static constexpr double eccentricity_squared()
    {
        return 1.0 -
            (Derived::b_m * Derived::b_m) /
            (Derived::a_m * Derived::a_m);
    }
};
\end{lstlisting}

\subsection{Concrete Policy}

A concrete policy supplies actual engineering content.

Example:

\begin{lstlisting}[language=C++]
struct Wgs84 : PlanetPolicyBase<Wgs84>
{
    using InertialFrame = frames::EarthCenteredInertial;
    using FixedFrame = frames::EarthCenteredEarthFixed;

    static inline constexpr double mu_m3_s2 = 3.986004418e14;
    static inline constexpr double omega_rad_s = 7.292115e-5;
    static inline constexpr double a_m = 6378137.0;
    static inline constexpr double b_m = 6356752.314245179;
    static inline constexpr double J2 = 1.08262668e-3;
};
\end{lstlisting}

\subsection{Runtime Algorithm}

A runtime algorithm operates on runtime data: state, covariance, measurements,
sensor samples, logs, buffers, and configuration values. It may be templated on
policies, but it should not own unrelated domain details.

\section{Design Drivers}

The design drivers are:

\begin{itemize}
    \item Zero runtime overhead
    \item Embedded suitability
    \item Compile-time validation
    \item Clear compiler diagnostics
    \item Avoidance of runtime polymorphism for static engineering choices
    \item Separation of domain definitions from runtime algorithms
    \item Ability to compose compatible policies
    \item Avoidance of inheritance hierarchy explosion
\end{itemize}

\section{Concepts as Capabilities}

Concepts in NavKit should generally represent capabilities.

For example, rather than creating a class hierarchy such as:

\begin{verbatim}
Planet
  -> RotatingPlanet
       -> EllipsoidPlanet
            -> J2Planet
\end{verbatim}

NavKit defines independent capability concepts:

\begin{lstlisting}[language=C++]
template<typename T>
concept PlanetPolicy = requires
{
    T::mu_m3_s2;
    typename T::InertialFrame;
    typename T::FixedFrame;
};

template<typename T>
concept RotatingPlanetPolicy =
    PlanetPolicy<T> &&
    requires { T::omega_rad_s; };

template<typename T>
concept EllipsoidPlanetPolicy =
    PlanetPolicy<T> &&
    requires
    {
        T::a_m;
        T::b_m;
    };

template<typename T>
concept J2PlanetPolicy =
    PlanetPolicy<T> &&
    requires
    {
        T::J2;
        T::a_m;
    };
\end{lstlisting}

This avoids combinatorial inheritance problems. A concrete policy simply
aggregates the facts it supports. The compiler determines which capabilities it
satisfies.

\section{Policy Aggregation}

Concrete policies should generally be flat aggregations of the facts and types
they support.

For example, \texttt{Wgs84} may simultaneously satisfy:

\begin{itemize}
    \item \texttt{PlanetPolicy}
    \item \texttt{RotatingPlanetPolicy}
    \item \texttt{EllipsoidPlanetPolicy}
    \item \texttt{J2PlanetPolicy}
\end{itemize}

without explicitly declaring any of those concepts.

This style is intentionally different from traditional object-oriented
inheritance. Algorithms constrain on the capabilities they need, not on a
particular place in a class hierarchy.

\section{Algorithm Constraints}

Runtime algorithms and higher-level policies should constrain template
parameters using the weakest concept that expresses their actual requirements.

Examples:

\begin{lstlisting}[language=C++]
template<PlanetPolicy Planet>
struct SphericalGravity;

template<J2PlanetPolicy Planet>
struct J2Gravity;

template<RotatingPlanetPolicy Planet, GravityPolicy Gravity>
struct MechanizationPcpf;
\end{lstlisting}

This means a planet that supports only \texttt{mu\_m3\_s2} may still be used
with spherical gravity, while a high-fidelity gravity policy may require
additional capabilities.

\section{CRTP Policy Bases}

CRTP policy bases are used for shared implementation, not for concept
satisfaction by declaration.

For example, \texttt{Wgs84} may inherit from
\texttt{PlanetPolicyBase<Wgs84>} to inherit helper functions. However, the
compiler determines concept satisfaction from the public interface, not from an
explicit \texttt{implements} declaration.

This distinction is important:

\begin{description}
    \item[Concepts] answer: ``Can this type be used here?''
    \item[CRTP bases] answer: ``What shared implementation can this family of
    policies reuse?''
\end{description}

\section{Dependency Containment}

Template types may indirectly encode lower-level policies. This is acceptable.

Example:

\begin{lstlisting}[language=C++]
using Planet = environment::Wgs84;
using Gravity = environment::J2<Planet>;
using Propagation = nav::MechanizationPcpf<Planet, Gravity>;
using Nav = Navigator<Filter, Sensors, Propagation, UpdatePolicy>;
\end{lstlisting}

The concrete \texttt{Navigator} type indirectly contains the planet and gravity
choice through its propagation policy. The important point is that
\texttt{Navigator} is semantically constrained on \texttt{PropagationPolicy},
not directly on planet, gravity, and frame policies.

\section{Naming Conventions}

NavKit shall use the following naming conventions for compile-time policies:

\begin{itemize}
    \item Concepts end with \texttt{Policy}
    \item CRTP bases end with \texttt{PolicyBase}
    \item Concrete policies use descriptive domain names
    \item Runtime algorithms use algorithmic names
\end{itemize}

Examples:

\begin{center}
\begin{tabular}{lll}
\toprule
\textbf{Role} & \textbf{Example} & \textbf{Meaning} \\
\midrule
Concept & \texttt{PlanetPolicy} & Required planet capability \\
Concept & \texttt{J2PlanetPolicy} & Planet with J2 support \\
CRTP Base & \texttt{PlanetPolicyBase} & Shared planet helpers \\
Concrete Policy & \texttt{Wgs84} & Specific planet definition \\
Concrete Policy & \texttt{J2<Planet>} & Specific gravity definition \\
Runtime Algorithm & \texttt{Navigator} & Orchestration object \\
\bottomrule
\end{tabular}
\end{center}

\section{Alternatives Considered}

\subsection{Runtime Polymorphism}

Runtime inheritance and virtual dispatch were rejected for static engineering
choices because they introduce runtime overhead and complicate embedded
deployment.

\subsection{Traditional Inheritance Hierarchies}

Deep inheritance hierarchies were rejected because capability combinations
quickly become difficult to represent. Concepts allow independent capabilities
without forcing a rigid taxonomy.

\subsection{Plain Templates Without Concepts}

Plain templates provide zero-overhead abstraction but produce poor compiler
diagnostics when requirements are not met.

\subsection{CRTP Without Concepts}

CRTP provides shared implementation but does not clearly constrain algorithm
interfaces at template boundaries.

\section{Consequences}

\subsection{Benefits}

\begin{itemize}
    \item Static configuration with no runtime dispatch
    \item Better compiler diagnostics
    \item Capability-based composition
    \item Avoidance of inheritance hierarchy explosion
    \item Reusable shared implementation
    \item Stronger architectural consistency
\end{itemize}

\subsection{Costs}

\begin{itemize}
    \item Increased template complexity
    \item Requires familiarity with C++20 concepts
    \item Potentially longer compile times
    \item Requires discipline to avoid propagating low-level policies too high
    in the architecture
\end{itemize}

\section{Guiding Rule}

Introduce a policy when an engineering assumption is:

\begin{itemize}
    \item fundamental to the algorithm,
    \item naturally static for a build or scenario,
    \item likely to have multiple valid implementations or capability levels,
    \item useful to express in the type system.
\end{itemize}

Use runtime objects for values that change during execution or represent
state, measurements, logs, buffers, and runtime configuration.

\section{Decision Summary}

NavKit shall use Concepts to express compile-time capabilities, CRTP policy
bases to share implementation, concrete policies to aggregate domain facts or
behavior, and runtime algorithms to operate on runtime data.

Concepts in NavKit are capability predicates, not base classes. Concrete
policies do not explicitly declare which concepts they implement. Instead,
algorithms constrain template parameters using the weakest concept that
expresses their actual requirements.

\end{document}
